% Pista pro-24i-q4 · Probabilitat
% Procedència: PAU Catalunya 2024 · Convocatòria de juny (incidències) · Sèrie 5 · Qüestió 4
\documentclass[11pt,a4paper]{article}
\usepackage{pista}
\pistainfo{Catalunya 2024}{Convocatòria de juny (incidències)}{Sèrie 5}

\begin{document}

\section*{\textcolor{blauPista}{Qüestió 4}}

\noindent La Rut fa servir el mètode següent per a fer els problemes de matemàtiques: tira un dau equilibrat i, si el resultat és com a màxim 4, pensa i resol el problema ella mateixa; si el resultat és 5 o 6, busca la solució del problema per Internet i la copia. Quan és ella qui ha pensat la solució, la resposta és correcta en el $75\,\%$ dels casos; quan copia la solució d'Internet, la resposta és correcta només en el $40\,\%$ dels casos.

\subsection*{\textit{a)} \normalfont Quina és la probabilitat que la solució d'un problema respost seguint aquest mètode sigui correcta? \hfill\textnormal{[0,75~punts]}}

\pista{Un diagrama d'arbre t'ajudarà a visualitzar els dos camins que porten a una resposta correcta: o bé la Rut ha pensat la solució i és correcta, o bé l'ha copiat i és correcta. Aplica el teorema de la probabilitat total. Ves amb compte amb les probabilitats inicials: com que el dau és equilibrat, l'esdeveniment ``com a màxim 4'' té probabilitat $\dfrac{4}{6}$ i, en canvi, l'esdeveniment ``5 o 6'' té probabilitat $\dfrac{2}{6}$.}

\subsection*{\textit{b)} \normalfont Quina és la probabilitat que un problema l'hagi resolt la Rut si sabem que la solució és correcta? \hfill\textnormal{[0,75~punts]}}

\pista{Aquest és un exemple clar on has d'aplicar el Teorema de Bayes, perquè et demanen una probabilitat condicionada en sentit \emph{invers} respecte de la informació que et donen: tu coneixes la probabilitat de ``la solució és correcta \emph{sabent que} l'ha resolt ella'', però et demanen la probabilitat de ``l'ha resolt ella \emph{sabent que} la solució és correcta''.}

\subsection*{\textit{c)} \normalfont Demà la Rut ha d'entregar 5 problemes de matemàtiques. Quina és la probabilitat que n'hi hagi almenys 4 de correctes? \hfill\textnormal{[1~punt]}}

\pista{Aquí tens una variable binomial $X \sim B(5, p)$, on $p$ és la probabilitat que un problema qualsevol sigui correcte (la que has calculat a l'apartat a). L'esdeveniment ``almenys 4 correctes'' equival a ``exactament 4 o exactament 5'': calcula $P(X = 4)$, calcula $P(X = 5)$ i suma els dos resultats.}

\end{document}
